\documentclass[11pt]{article}
\usepackage[margin=1.1in]{geometry}
\usepackage[T1]{fontenc}
\usepackage{amsmath,amssymb,amsthm}
\usepackage{booktabs}
\usepackage[hidelinks]{hyperref}

\newtheorem{theorem}{Theorem}[section]
\newtheorem{proposition}[theorem]{Proposition}
\newtheorem{lemma}[theorem]{Lemma}
\newtheorem{corollary}[theorem]{Corollary}
\newtheorem{conjecture}[theorem]{Conjecture}
\theoremstyle{definition}
\newtheorem{definition}[theorem]{Definition}
\newtheorem{example}[theorem]{Example}
\theoremstyle{remark}
\newtheorem{remark}[theorem]{Remark}

\title{A Vertex-Gap Obstruction for Low-Degree Strip Pairs\\ in the Plane Jacobian Conjecture}
\author{Dan Clemens Posch\thanks{Email: \texttt{dc@dcpos.ch}. This work was produced in close collaboration with Claude (Fable, an AI system by Anthropic), which developed much of the mathematics and all of the software under the author's direction and review. The author takes responsibility for every claim. The division of labor and the adversarial verification process are described in Section \ref{sec:computation}.}}
\date{August 2026}

\begin{document}
\maketitle

\begin{abstract}
Below degree $125$, the reduction program of Guccione, Guccione, Valqui and Horruitiner leaves a single open family of possible counterexample degrees to the plane Jacobian conjecture, $(\deg P,\deg Q)=(108,72)$. One of its two reduced subcases is a pair of parallel Newton-polygon strips with $[P,Q]=x^2$. We develop an obstruction theory for such strip pairs: a vertex normalization at a Minkowski vertex, a gap condition on the bottom corners, and a counting argument on a self-contained near-origin block. At $(k,d_2)=(2,2)$ this empties the generic chart; the $(8,28)$ family is the explicit ten-step case. Our main theorem describes the block at every cell $k,d_2\geq 2$ in characteristic zero: the depth-two block variety is $\{A\ \text{binomial}\}\cap(\{a_2=0\}\cup\{R_{k,d_2}=0\})$, where $A$ is the low $P$-column and $R_{k,d_2}$ is an explicit logarithmic-residue functional whose one-term case is the $(2,2)$ theorem. The rigidity input is an elementary lemma: if $AC'-\nu A'C$ is a nonzero constant, with $A,C$ polynomials and $\nu\geq1$ an integer, then $\deg A\leq 1$. Deeper Minkowski columns, the $k=1$ cells, and non-strip shapes remain open. All claims are machine-verified in exact arithmetic.
\end{abstract}

\section{Introduction}\label{sec:intro}

The Jacobian conjecture of Keller asserts that a polynomial map $F\colon \mathbb{C}^n\to\mathbb{C}^n$ with $\det JF$ a nonzero constant is invertible. In July 2026 Alp\"oge announced an explicit cubic counterexample in dimension three \cite{Alpoge}; adding dummy variables, the conjecture fails for all $n\geq 3$. The plane case $\mathrm{JC}_2$ remains open and has independent structure: graded-equivariant Keller maps of $\mathbb{C}^2$, for example, are automorphisms for every weight pattern \cite{Shaska}.

Degree is the oldest invariant. Moh proved that a plane counterexample $(P,Q)$ must have $\max(\deg P,\deg Q)>100$ \cite{Moh}; Heitmann described the possible degree pairs and Newton-polygon corners \cite{Heitmann}. \.Zo{\l}\k{a}dek \cite{Zoladek} and, independently, Sigray \cite{Sigray} bounded the topological degree below by $6$. The sharpest current information is the program of Guccione, Guccione and Valqui (GGV): a discrete-geometry analysis of the Newton polygons of a hypothetical counterexample \cite{GGV1,GGV2,GGV3}, made algorithmic in \cite{GGV4}, tabulating $34$ admissible corner families with $\max(\deg P,\deg Q)\leq 150$. With Horruitiner they raised the degree bound from $100$ to $108$ \cite{GGHV}. Below $125$ one family survives: corner $A_0=(8,28)$, degrees $(\deg P,\deg Q)=(108,72)$.

Proposition 4.3 of \cite{GGHV} reduces this family to two subcases: find $(P,Q)$ with $[P,Q]:=P_xQ_y-P_yQ_x=x^2$, with $N(P)$ and $N(Q)$ in prescribed small polygons and prescribed vertices attained. In subcase (2) both polygons are narrow lattice \emph{strips} along $(1,2)$. A computational campaign discarded that subcase over characteristic zero (conditional on Proposition 4.3; Section \ref{sec:computation}); the generic chart empties by a short substitution chain ending in $-1=0$. This note isolates that mechanism, the \emph{vertex-gap obstruction}, and develops it as far as it provably goes.

For reduced strip pairs with $[P,Q]=x^k$ and direction $(1,d_2)$ (Section \ref{sec:setup}), we prove two counting propositions that locate a self-contained near-origin block of bracket equations (Section \ref{sec:gap}); the vertex-gap theorem at $(k,d_2)=(2,2)$, emptying the chart $a_2a_6\neq 0$ by a surplus relation of pinning type in the radical of the block ideal (Section \ref{sec:gap}); variants at $(2,3)$ and $(2,4)$ (support rigidity of a $P$-column plus a residue hyperplane), rigidity certificates at $(3,3)$, $(4,3)$, $(5,3)$, and the absence of any near-origin obstruction for $k\geq 3$ at $d_2=2$ (Section \ref{sec:variants}); and the explicit ten-event collapse for the $(8,28)$ generic chart (Section \ref{sec:example}). Section \ref{sec:functional} proves the main structural result: a logarithmic-residue functional $R_{k,d_2}$, whose one-term case is the $(2,2)$ theorem, describes the depth-two block variety at every cell $k,d_2\geq 2$ in characteristic zero, subsuming and completing the per-cell results of Section \ref{sec:variants}. Section \ref{sec:coverage} records which GGV families the theory covers; Section \ref{sec:computation} describes verification and AI assistance.

\section{Setup: strip pairs, valuations, and the vertex normalization}\label{sec:setup}

Throughout, $K$ is a field, $P,Q\in K[x,y]$, and $[P,Q]:=P_xQ_y-P_yQ_x$. We write $P=\sum_p a_p x^{p_1}y^{p_2}$ and $Q=\sum_q b_q x^{q_1}y^{q_2}$ over lattice points $p=(p_1,p_2)$, $q=(q_1,q_2)\in\mathbb{Z}_{\geq 0}^2$, and $N(P)$ for the Newton polygon (the convex hull of the support). The whole theory rests on one formula: since
\[
[x^{p_1}y^{p_2},\,x^{q_1}y^{q_2}] \;=\; \det(p,q)\, x^{p_1+q_1-1}y^{p_2+q_2-1},
\qquad \det(p,q):=p_1q_2-p_2q_1,
\]
the coefficient of $[P,Q]$ attached to a \emph{Minkowski point} $m\in\mathbb{Z}^2$ is
\begin{equation}\label{eq:key}
E_m \;:=\; \sum_{p+q=m} \det(p,q)\, a_p b_q ,
\end{equation}
the coefficient of the monomial $x^{m_1-1}y^{m_2-1}$. The equation $[P,Q]=x^k$ says: $E_m=0$ for every $m$ in the Minkowski sum $N(P)+N(Q)$ except $E_{(k+1,1)}=1$, since $x^k$ sits at the Minkowski point $(k+1,1)$. We call the $E_m$ \emph{keys} and index them by $m$.

Fix a primitive direction $d=(1,d_2)$ with $d_2\geq 1$ and define the weight and the level of a lattice point by
\[
w(i,j) \;:=\; d_2 i - j, \qquad \ell(i,j) \;:=\; d_2 i + j .
\]
The weight $w$ is constant on lines parallel to $d$ and vanishes on the line $\mathbb{R}d$ through the origin; the level $\ell$ is injective on each vertical column and orders a strip from the origin outward. Both gradings are additive: $w(p+q)=w(p)+w(q)$, likewise for $\ell$. If $w(p)=w(q)=0$ then $p,q$ are proportional and $\det(p,q)=0$: the extreme $w$-stratum of the bracket vanishes identically.

\begin{definition}\label{def:strip}
A \emph{reduced strip pair} of type $(k,d_2)$ consists of $P,Q$ with $[P,Q]=x^k$, $k\geq 1$, such that:
\begin{enumerate}
\item[(S1)] $N(P)$ and $N(Q)$ are \emph{strips} along $d=(1,d_2)$: every support point $p$ of $P$ satisfies $0\leq w(p)\leq w_P$, and every support point $q$ of $Q$ satisfies $0\leq w(q)\leq w_Q$, for integers $w_P,w_Q\geq 1$; the top edges ($w=0$) of both polygons lie on the line $\mathbb{R}d$ through the origin.
\item[(S2)] The bottom edges are parallel to $d$ and start at \emph{saturated} bottom corners $p_0$, $q_0$: these are vertices of the polygons, $w(p_0)=w_P$, $w(q_0)=w_Q$, and $a_{p_0}\neq 0$, $b_{q_0}\neq 0$.
\item[(S3)] (Block containment.) $P$ may have support at every lattice point of its strip in columns $x=1,2$, and $Q$ in columns $x=k,k+1$; the strips extend at least that far. No support on the $y$-axis except possibly the origin.
\end{enumerate}
\end{definition}

The reader should keep in mind the example that motivated everything, the $(8,28)$ subcase (2) of \cite{GGHV}: $N(P)=\mathrm{hull}\{(0,0),(1,0),(8,14),(8,16)\}$ and $N(Q)=\mathrm{hull}\{(0,0),(2,1),(12,21),(12,24)\}$ with $[P,Q]=x^2$. Here $d=(1,2)$, $w(i,j)=2i-j$, $P$ occupies weights $\{0,1,2\}$ and $Q$ occupies $\{0,1,2,3\}$, the top edges $(0,0)$ to $(8,16)$ and $(0,0)$ to $(12,24)$ lie on $\mathbb{R}d$, and the bottom corners are $p_0=(1,0)$, $q_0=(2,1)$, forced nonzero because Proposition 4.3 prescribes vertex attainment.

\medskip
\noindent\textbf{The vertex normalization.}
The right-hand side $x^k$ lives at the Minkowski point $(k+1,1)$. For a strip pair the natural position for this point is the bottom vertex $p_0+q_0$ of $N(P)+N(Q)$, the lower endpoint of the long Minkowski edge of direction $d$ obtained by concatenating the two parallel bottom edges. In the motivating example $(3,1)=(1,0)+(2,1)$ is such a vertex. We take this as a hypothesis:

\begin{itemize}
\item[(i)] (\emph{vertex normalization}) $(k+1,1)=p_0+q_0$, this is the unique decomposition of $(k+1,1)$ as a sum of a support point of $P$ and one of $Q$, and $\det(p_0,q_0)=\pm 1$.
\end{itemize}

Under (i) the key at the vertex reads $E_{(k+1,1)}=\det(p_0,q_0)\,a_{p_0}b_{q_0}=1$, so both bottom-corner coefficients are units. We call this the \emph{vertex equation}: the valuation of $x^k$ selects the bottom vertex of the tower of parallel edges, and normalizes the corner product. Hypothesis (i) is very rigid:

\begin{lemma}\label{lem:enum}
Under (i), $\{p_0,q_0\}=\{(1,0),(k,1)\}$.
\end{lemma}

\begin{proof}
Write $p_0=(\alpha,\beta)$, $q_0=(k+1-\alpha,1-\beta)$ with all entries $\geq 0$, so $\beta\in\{0,1\}$. If $\beta=0$ then $\det(p_0,q_0)=\alpha=\pm1$ forces $\alpha=1$, giving $p_0=(1,0)$, $q_0=(k,1)$. If $\beta=1$ then $\det(p_0,q_0)=-(k+1-\alpha)=\pm 1$ forces $\alpha=k$, giving the same pair with the roles swapped.
\end{proof}

By Lemma \ref{lem:enum} we may and do \emph{normalize the labels} so that $p_0=(1,0)$ and $q_0=(k,1)$; if the corner $(k,1)$ belongs to $P$ instead, replace $(P,Q)$ by $(Q,-P)$, which has the same bracket and swaps the roles. The widths are then determined: $w_P=w(1,0)=d_2$ and $w_Q=w(k,1)=kd_2-1$. Each full strip column is a ladder of consecutive lattice points: column $x=c$ of $P$ is $\{(c,j): d_2c-w_P\leq j\leq d_2c\}$, and likewise for $Q$, so at type $(k,d_2)$ a full $P$-column has $d_2+1$ points and a full $Q$-column has $kd_2$ points.

\medskip
\noindent\textbf{The gap condition.} The second hypothesis is a single inequality:
\begin{itemize}
\item[(ii)] (\emph{gap condition}) $k\geq 2$; equivalently, by Lemma \ref{lem:enum}, $\max((p_0)_x,(q_0)_x)\geq 2$.
\end{itemize}
When $k\geq 2$, the polygon $Q$ has support columns $x=1,\dots,k-1$ strictly between the $y$-axis and the column of its bottom corner. Because the bottom edge of $N(Q)$ runs from $(0,0)$ to $(k,1)$, the points of such a \emph{gap column} $x=g$ are $(g,y)$ with $1\leq y\leq d_2g$, strictly above the bottom edge. These columns produce bracket keys below and to the left of the vertex $(k+1,1)$, and Section \ref{sec:gap} shows that they vanish on the solution set. When $k=1$ there is no gap column on either side, and no such forcing exists. Under (i), ``$k\geq 2$'' is equivalent to the existence of a gap column on exactly one side.

\section{The gap condition and the counting theorem}\label{sec:gap}

We now assume (i) and (ii) in the normalized labels: $p_0=(1,0)$, $q_0=(k,1)$, $w_P=d_2$, $w_Q=kd_2-1$, and $a_{p_0}$, $b_{q_0}$ units. The plan: first the gap columns die; then a small block of keys near the origin becomes a self-contained linear system in the nearby coefficients; counting its equations against its unknowns explains, and in favorable cells proves, the collapse.

\begin{lemma}[gap kill]\label{lem:gapkill}
On the solution set, every coefficient of $Q$ in the gap columns $x=1,\dots,k-1$ vanishes.
\end{lemma}

\begin{proof}
Induct on the column $g$ and, inside a column, on the level $\ell$. Let $(g,y)$ be a gap point, so $1\leq y\leq d_2g$. The key at the sub-vertex Minkowski point $p_0+(g,y)=(1+g,y)$ contains the monomial $\det(p_0,(g,y))\,a_{p_0}b_{(g,y)}=y\,a_{p_0}b_{(g,y)}$ with $y\geq 1$. Any other pair $(p',q')$ with $p'+q'=(1+g,y)$ either has $q'$ the origin (then $\det(p',q')=0$) or has $p'$ in a $P$-column $x\geq 1$, hence $q'$ in a gap column $\leq g$ with either smaller column or smaller level than $(g,y)$; by induction $b_{q'}=0$. So the key reduces to a single monomial with unit cofactor $a_{p_0}$ and nonzero constant $y$, forcing $b_{(g,y)}=0$. (This uses that $y$ is invertible: $1\leq y\leq d_2(k-1)$, a first source of excluded small primes.)
\end{proof}

In the $(8,28)$ example the gap column is $x=1$ of $Q$, with points $(1,1)$, $(1,2)$; the sub-vertex keys $E_{(2,1)}=a_{p_0}b_{(1,1)}$ and $E_{(2,2)}=2a_{p_0}b_{(1,2)}$ kill it. The mechanism needs the absence of $y$-axis support (S3): a $P$-coefficient at $(0,a)$ would contribute $\det((0,a),(g',b))\neq 0$ terms to the same keys, and the forcing fails. That is why the pentagon-shaped subcase (1) of \cite{GGHV} is outside the theory (Section \ref{sec:coverage}).

\medskip
\noindent\textbf{The near-origin block.} Define the \emph{block} to be the set of keys in Minkowski columns $k+1$ and $k+2$. After the gap kill, a pair $(p,q)$ with $a_pb_q$ alive can land in the block only with $p$ in $P$-column $1$ or $2$ and $q$ in $Q$-column $k$ or $k+1$: the origin pairs have $\det=0$, gap columns are dead, and larger $P$-columns would push $q$ into a dead or negative column. Conversely such pairs land only in the block. So the block is \emph{self-contained}: its equations involve the coefficients in $P$-columns $1,2$ and $Q$-columns $k,k+1$ only, no matter how long the strips are. All counting can be done locally, and the block has fixed size depending only on $(k,d_2)$.

\begin{proposition}[key count]\label{prop:A}
For all $k,d_2\geq 1$, the number of forced-to-vanish block keys that are not identically zero is $n_{\mathrm{keys}}=2(w_P+w_Q)-1$.
\end{proposition}

\begin{proof}
Column $k+1$ of $N(P)+N(Q)$ is the ladder $(k+1,y)$, $1\leq y\leq d_2(k+1)$, i.e.\ strata $w=0,\dots,w_P+w_Q$; the vertex $(k+1,1)$ is the bottom stratum $w=w_P+w_Q$. The stratum $w=0$ is identically zero: $w(p)+w(q)=0$ with both weights $\geq 0$ forces $w(p)=w(q)=0$, so $p,q\in\mathbb{R}d$ and $\det(p,q)=0$. Every stratum $1\leq w\leq w_P+w_Q$ occurs after the gap kill: for $w\geq w_P$ take the pair $(p_0,q)$ with $q$ in the full $Q$-column $k$ at weight $w-w_P$, whose determinant is $q_2\geq 1$; for $1\leq w\leq w_P-1$ take $p=(1,d_2-w)$ in $P$-column $1$ and $q=(k,kd_2)$ the top of $Q$-column $k$, with $\det(p,q)=kw\neq 0$. Distinct pairs give distinct monomials $a_pb_q$, so no cancellation occurs and each such key is a genuine equation. Removing the vertex key leaves $w_P+w_Q-1$ keys. Column $k+2$ is the same ladder argument with bottom point $(k+2,1+d_2)$ on the long Minkowski edge (witnesses: $(p_0,q)$ with $q$ in column $k+1$, and $p=(1,d_2-w)$ against the top of column $k+1$, $\det=(k+1)w$); there is no vertex to remove, giving $w_P+w_Q$ keys. Total $2(w_P+w_Q)-1$.
\end{proof}

\begin{proposition}[eliminations and extra keys]\label{prop:B}
All $2w_Q+1$ non-corner coefficients $b_q$ with $q$ in $Q$-columns $k$ and $k+1$ are eliminable by unit pivots, each from its own key at $m=p_0+q$; this consumes the block keys of strata $w_P\leq w\leq w_P+w_Q$ (vertex excluded). The remaining \emph{extra keys} are the $2(w_P-1)$ keys of strata $1\leq w\leq w_P-1$ in the two block columns.
\end{proposition}

\begin{proof}
The key at $p_0+q$ contains $\det(p_0,q)\,a_{p_0}b_q=q_2\,a_{p_0}b_q$ with $q_2\geq 1$ and unit cofactor $a_{p_0}$, and this is the only monomial of that key containing $b_q$ (the pair with second factor $q$ is unique). The map $q\mapsto p_0+q$ is injective, hits the strata $w(p_0+q)=w_P+w(q)\in[w_P,w_P+w_Q]$ in each block column, and misses only the vertex ($q=q_0$). Triangularity: in any other monomial $(p',q')$ of the key at $p_0+q$ one has $\ell(p')>\ell(p_0)$ (a point of $P$-column $1$ other than $p_0$ has $\ell=d_2+j>d_2$, and column $2$ has $\ell\geq 2d_2$), hence $\ell(q')<\ell(q)$, so eliminating in increasing level of $q$ is well founded and each substitution inserts only already-expressed or lower-level variables. The counts follow: $2w_Q+1$ keys consumed (a full $Q$-column has $kd_2=w_Q+1$ points, two columns minus the corner), and $n_{\mathrm{keys}}-(2w_Q+1)=2(w_P-1)$ extra keys, namely the strata $1\leq w\leq w_P-1$.
\end{proof}

Propositions \ref{prop:A} and \ref{prop:B} hold for every $k,d_2\geq 1$, provided the integer determinants used as pivots are invertible. The content of the extra keys, after back-substituting the eliminated $b$'s, varies from cell to cell. An extra key is of \emph{pinning type} if its back-substituted form is $(\text{unit})\cdot M\cdot b_{q_0}$ with $M$ a nonzero monomial in the $P$-coefficients: such a key forces $b_{q_0}=0$ on the chart where $M$ is invertible, contradicting the vertex equation. The main theorem says this happens at $(k,d_2)=(2,2)$, with $M=a_2^2a_6$.

\begin{theorem}[vertex-gap obstruction, $(k,d_2)=(2,2)$]\label{thm:22}
Let $(P,Q)$ be a reduced strip pair of type $(2,2)$ satisfying \textup{(i)}, with strips of any lengths containing the block columns, over a field of characteristic not in $\{2,3,5\}$. Write
\[
a_1=a_{(1,0)},\; a_2=a_{(1,1)},\; a_3=a_{(1,2)},\; a_4=a_{(2,2)},\; a_5=a_{(2,3)},\; a_6=a_{(2,4)}
\]
for the block $P$-coefficients ($a_1$ a unit) and $b_3=b_{(2,1)}$ (a unit), $b_4=b_{(2,2)}$, $b_5=b_{(2,3)}$, $b_6=b_{(2,4)}$, $b_7=b_{(3,3)}$, $\dots$, $b_{10}=b_{(3,6)}$ for the block $Q$-coefficients. Then every solution satisfies $a_3=0$ and $a_2^2a_6=0$. Thus the locus $a_2a_6\neq 0$ is empty, and the block variety is contained in $\{a_3=0\}\cap(\{a_2=0\}\cup\{a_6=0\})$.
\end{theorem}

\begin{proof}
Here $w_P=2$, $w_Q=3$, $p_0=(1,0)$, $q_0=(2,1)$; by Proposition \ref{prop:A} there are $9$ block keys besides the vertex, and by Proposition \ref{prop:B} seven of them eliminate $b_4,\dots,b_{10}$. After the gap kill ($b_{(1,1)}=b_{(1,2)}=0$), the nine keys, in elimination order, reduce as follows (each line lists the Minkowski point, the key, and its consequence after substituting the previous lines):
\begin{align*}
(3,2)&: & 2a_1b_4-a_2b_3&=0 &&\Rightarrow\; b_4=\tfrac{a_2b_3}{2a_1}\\
(3,3)&: & 3a_1b_5-3a_3b_3&=0 &&\Rightarrow\; b_5=\tfrac{a_3b_3}{a_1}\\
(3,4)&: & 4a_1b_6+a_2b_5-2a_3b_4&=0 &&\Rightarrow\; b_6=0 \quad [\mathrm{C}1]\\
(3,5)&: & 2a_2b_6-a_3b_5&=0 &&\Rightarrow\; -\tfrac{a_3^2b_3}{a_1}=0 \;\Rightarrow\; a_3=0\\
(4,3)&: & 3a_1b_7-2a_4b_3&=0 &&\Rightarrow\; b_7=\tfrac{2a_4b_3}{3a_1}\\
(4,4)&: & 4a_1b_8-4a_5b_3&=0 &&\Rightarrow\; b_8=\tfrac{a_5b_3}{a_1}\\
(4,5)&: & 5a_1b_9+a_2b_8-3a_3b_7+2a_4b_5-2a_5b_4-6a_6b_3&=0 &&\Rightarrow\; b_9=\tfrac{6a_6b_3}{5a_1} \quad [\mathrm{C}2]\\
(4,6)&: & 6a_1b_{10}+2a_2b_9-2a_3b_8+4a_4b_6-4a_6b_4&=0 &&\Rightarrow\; b_{10}=-\tfrac{a_2a_6b_3}{15a_1^2}\\
(4,7)&: & 3a_2b_{10}-a_3b_9+2a_5b_6-2a_6b_5&=0 &&\Rightarrow\; -\tfrac{1}{5}\,\tfrac{a_2^2a_6b_3}{a_1^2}=0 .
\end{align*}
Two cancellations create the surplus. In $(3,4)$, substituting $b_5$ and $b_4$ gives $a_2b_5-2a_3b_4=\frac{a_2a_3b_3}{a_1}-\frac{a_2a_3b_3}{a_1}=0$, so $b_6=0$ [C1]. In $(4,5)$, the pairs $-3a_3b_7+2a_4b_5$ and $a_2b_8-2a_5b_4$ cancel identically (both differences vanish even before $a_3=0$), leaving $b_9$ proportional to $a_6$ alone [C2]. The key $(3,5)$ then carries a single monomial $-a_3^2b_3/a_1$ with unit cofactors, forcing $a_3=0$ on the solution set. The final line is the surplus key: nine equations against eight eliminable unknowns ($b_4,\dots,b_{10}$ and $a_3$). Its back-substituted form is of pinning type with $M=a_2^2a_6$, so $a_2^2a_6\,b_3=0$; since $b_3$ is a unit by the vertex equation, $a_2^2a_6=0$. The pivot constants used are $\{2,3,4,3,4,5,6\}$, the gap pivots $\{1,2\}$, and the final $1/5$: characteristic outside $\{2,3,5\}$ suffices.
\end{proof}

\begin{remark}[radical, not ideal, membership]\label{rem:radical}
The step $a_3=0$ is extracted from the relation $a_3^2b_3/a_1=0$, which is quadratic in $a_3$. The chart-free form of the conclusion is therefore a \emph{radical} membership: $u:=a_2^2a_6b_{q_0}$ satisfies $u^2\in I$, where $I$ is the ideal generated by the block, gap and vertex equations, but $u\notin I$. A witness over the dual numbers $K[\varepsilon]/(\varepsilon^2)$: set $a_3=\varepsilon$, $a_2=1$, $a_6=\varepsilon c$; all generators of $I$ can be satisfied while $u=\varepsilon c\, b_{q_0}\neq 0$. Every variety-level statement in Theorem \ref{thm:22} is unaffected.
\end{remark}

\begin{remark}[position of the surplus key]\label{rem:position}
In the general labels, the surplus key is the topmost non-vanishing key of the \emph{outer} block column: stratum $w=1$ of Minkowski column $k+2$, i.e.
\[
m_{\mathrm{surplus}}=(k+2,\; d_2(k+2)-1),
\]
which at $(2,2)$ is $(4,7)$. The $w=1$ key of the inner column, $(k+1,d_2(k+1)-1)=(3,5)$, is the rigidity event that forces $a_3=0$, not the surplus. A bare count of equations minus unknowns equals $1$ even in cells where no obstruction exists (Section \ref{sec:variants}); it is the pinning-type \emph{content} of the surplus key, not the count, that discriminates.
\end{remark}

\begin{remark}[excluded primes]
The set $\{2,3,5\}$ is the set of primes dividing some constant used by this derivation, and each occurs as a pivot, so the set is sharp for the proof. Exact rank computations at random points suggest the obstruction persists modulo $2$, $3$ and $5$ as well; we make no claim in those characteristics.
\end{remark}

\section{Variant exponents and an impossibility result}\label{sec:variants}

Outside $(2,2)$, by Proposition \ref{prop:B} the block has $2(w_P-1)=2(d_2-1)$ extra keys, so the single-surplus picture is special to $d_2=2$. To compute the extra keys uniformly we assemble each strip column into a generating polynomial. Write
\[
A(y)=\sum_{i=0}^{d_2} a_{(1,i)}y^i,\quad
B(y)=\sum_{j=d_2}^{2d_2} \beta_j y^j\ (\beta_j:=a_{(2,j)}),\quad
C(y)=\sum_{j=1}^{kd_2} b_{(k,j)}y^j,\quad
E(y)=\sum_{j=d_2+1}^{(k+1)d_2} b_{(k+1,j)}y^j
\]
for $P$-columns $1,2$ and $Q$-columns $k,k+1$. From $\det((1,i),(k,j))=j-ik$ and $\det((2,i),(k,j))=2j-ik$ one checks directly that the block keys are the coefficients of
\begin{equation}\label{eq:odes}
\text{column } k{+}1:\;\; AC'-kA'C, \qquad
\text{column } k{+}2:\;\; AE'-(k{+}1)A'E+2BC'-kB'C,
\end{equation}
the key at Minkowski point $(k{+}1,s)$ resp.\ $(k{+}2,s)$ being the coefficient of $y^{s-1}$, and the vertex equation the constant term $a_{p_0}b_{q_0}$ of the first expression. This assembly is machine-checked against an independent lattice-determinant enumeration of every block key. Proposition \ref{prop:B} solves the lower coefficients of these two linear ODE-like relations for $C$ and $E$; the extra keys are the would-be next coefficients of the solved series. We normalize the units $a_{p_0}=b_{q_0}=1$ in this section; unit factors are restored routinely.

\subsection{Type $(2,3)$: support rigidity plus one hyperplane}

At $(k,d_2)=(2,3)$ we have $A=1+a_2y+a_3y^2+a_4y^3$, $B=\beta_3y^3+\dots+\beta_6y^6$, and there are four extra keys: the inner pair $L_1,L_2$ (Minkowski $(3,7),(3,8)$, the coefficients of $y^6,y^7$ in $AC'-2A'C$) and the outer pair $L_3,L_4$ ($(4,10),(4,11)$). Back-substitution gives the exact forms
\begin{align*}
10L_1&=2a_3^3-4a_2^2a_3^2+2a_2^3a_4+19a_2a_3a_4-25a_4^2,\\
10L_2&=-2a_2a_3^3+a_2^2a_3a_4+8a_3^2a_4+a_2a_4^2,
\end{align*}
while $L_3,L_4$ are longer (18 monomials each), linear in $\beta_3,\dots,\beta_6$. The forms $L_1,L_2$ involve only the $P$-column $A$: no pinning of $b_{q_0}$ is available here, and on the computed grid no extra key at $d_2\geq 3$ reduces to a single monomial. The inner pair forces a support degeneration instead.

\begin{theorem}[rigidity at $(2,3)$]\label{thm:23}
Over a field of characteristic not in $\{2,3,5\}$, $V(L_1,L_2)=\{a_3=a_4=0\}$: the inner extra keys force $A=1+a_2y$, a binomial.
\end{theorem}

\begin{proof}[Proof 1 (resultant certificate)]
$\operatorname{Res}_{a_4}(10L_1,10L_2)=-3200\,a_3^7$, and the leading coefficient of $10L_1$ in $a_4$ is the constant $-25$, so for every fixed $(a_2,a_3)$ the two polynomials have a common root in $a_4$ only if the resultant vanishes: any common zero has $a_3=0$. On that slice $10L_2=a_2a_4^2$ and $10L_1=2a_2^3a_4-25a_4^2$; if $a_4\neq0$ the first forces $a_2=0$ and then the second reads $-25a_4^2\neq0$. So $a_4=0$. Conversely $a_3=a_4=0$ kills every monomial.
\end{proof}

\begin{proof}[Proof 2 (structural)]
The inner extras vanish when and only when the column identity $C'A-2A'C=1$ holds as polynomials, with $C$ the solved column of degree $\leq 6$. At a root $r$ of $A$ of multiplicity $\geq2$ the left side is divisible by $y-r$ while the right side is $1$: so $A$ is squarefree, and at each simple root $-2A'(r)C(r)=1$, whence $A'(r),C(r)\neq0$. Differentiating the identity twice and evaluating at $r$ yields $3A''(r)^2=A'(r)A'''(r)$, i.e.\ $A$ divides $N(A):=3A''^2-A'A'''$. If $\deg A=3$ then $N=(12a_3^2-6a_2a_4)+60a_3a_4y+90a_4^2y^2$ has degree $<3$, so divisibility forces $N=0$, but its top coefficient $90a_4^2$ is nonzero. If $\deg A=2$ then $N=12a_3^2$ is a nonzero constant. Hence $\deg A\leq 1$.
\end{proof}

\begin{theorem}[complete block variety at $(2,3)$]\label{thm:23full}
On the locus $a_3=a_4=0$ one has $L_4\equiv 0$ identically and $L_3=\frac{a_2^3}{21}\,R$, where
\[
R \;=\; a_2^2\beta_4-5a_2\beta_5+15\beta_6 .
\]
Hence $V(L_1,L_2,L_3,L_4)=\{a_3=a_4=0,\ a_2=0\}\ \cup\ \{a_3=a_4=0,\ R=0\}$, so the block variety contains no point with all block coordinates nonzero. The inner (rigidity) part uses only characteristics outside $\{2,3,5\}$; the outer computation excludes primes in $\{2,3,5,7,11\}$.
\end{theorem}

The coefficients $1,-5,15$ are $(-1)^j\binom{j}{4}$ for $j=4,5,6$; Section \ref{sec:functional} explains where they come from. Two contrasts with Theorem \ref{thm:22}: the forced content here is \emph{support rigidity} ($a_3=a_4=0$ are interior coefficients, so hypothesis (i) is untouched and no contradiction with the vertex equation arises), and there is no pinning ($b_{q_0}$ remains free, and on the main component $a_2\neq0$ the residual condition $R=0$ is a three-term linear form in $\beta_4,\beta_5,\beta_6$, satisfiable with every coordinate nonzero). The $(2,2)$ single-monomial collapse is special.

\subsection{Type $(2,4)$, and certificates at $(3,3)$, $(4,3)$, $(5,3)$}

\begin{theorem}[rigidity and block variety at $(2,4)$]\label{thm:24}
Let $(k,d_2)=(2,4)$, $A=1+py+qy^2+ry^3+sy^4$, over characteristic $0$ (the proof excludes only primes in $\{2,3,5,7\}$ for the rigidity part, plus $11$ for the outer part). The three inner extras force $q=r=s=0$, so $A$ is again a binomial; on that locus two of the three outer extras vanish identically and the third equals $-\frac{p^4}{55}R_4$ with
\[
R_4=p^4\beta_4-5p^3\beta_5+15p^2\beta_6-35p\beta_7+70\beta_8 .
\]
Hence $V(\text{all six extras})=\{q=r=s=0,\ p=0\}\cup\{q=r=s=0,\ R_4=0\}$, and again no point has all block coordinates nonzero.
\end{theorem}

\begin{proof}[Proof of the rigidity part]
As in Theorem \ref{thm:23}, the inner extras vanish iff $C'A-2A'C=1$, forcing $A$ squarefree and $A\mid N(A)=3A''^2-A'A'''$. Compute $N-336sA=E_0+60E_1y+30E_2y^2$ (the top two coefficients cancel), where $E_0=12q^2-6pr-336s$, $E_1=qr-6ps$, $E_2=3r^2-8qs$. If $\deg A=4$ ($s\neq0$), divisibility of the degree-$\leq2$ remainder by $A$ forces $E_0=E_1=E_2=0$; the identities
\[
9sE_0=84s(q^2-36s)-3qE_2+9rE_1,\qquad 27r^2-2q^3=9E_2-2q(q^2-36s)
\]
then give $q^2=36s$ and $27r^2=2q^3$, and $E_1=0$ with $q\neq0$ gives $r=pq/6$; these force $(p,q,r,s)=(4u,6u^2,4u^3,u^4)$, i.e.\ $A=(1+uy)^4$ with $u\neq0$, which is not squarefree: contradiction. The cases $\deg A=3,2$ are as before. Hence $\deg A\leq1$.
\end{proof}

\begin{remark}
Every perfect power $(1+uy)^\delta$ with $\delta\geq4$ satisfies the divisibility $A\mid N(A)$; it is squarefreeness that eliminates it. Any direct extension of this argument to general $d_2$ must keep both halves; the proof in Section \ref{sec:functional} instead bypasses the divisibility entirely (Theorem \ref{thm:ode}).
\end{remark}

\begin{theorem}[rigidity certificates at $(k,3)$, $k=3,4,5$]\label{thm:k3}
At $(k,d_2)=(3,3),(4,3),(5,3)$ the inner extra keys again force $a_3=a_4=0$ ($A$ binomial); consequently these block varieties contain no point with all block coordinates nonzero. The proof is a machine-verified pair of exact Sylvester resultant computations for each cell: an $a_4$-resultant of two inner extras equal to $c\,a_3^m$ with constant $a_4$-leading coefficient, followed by an $a_2$-resultant on the slice $a_3=0$ equal to $c'a_4^N$; the exponents $(m,N)$ are $(15,12)$, $(26,34)$, $(40,44)$ respectively. The certificates are exact rational arithmetic, locked by executable assertions in the repository (Section \ref{sec:computation}).
\end{theorem}

Theorem \ref{thm:k3} was certified computation without a structural derivation; Theorem \ref{thm:R} now provides one at every cell $k,d_2\geq2$, including the then-open $(3,4)$, $(4,4)$, $(5,4)$ and all $d_2\geq5$, and the certificates remain as independent corroboration.

\subsection{Impossibility for $k\geq 3$ at $d_2=2$}

The gap condition (ii) holds for every $k\geq2$, but the $(2,2)$ collapse does not generalize in $k$.

\begin{theorem}[no near-origin obstruction, $k\geq3$, $d_2=2$]\label{thm:kbig}
Fix $k\geq3$ and $d_2=2$, over a field of characteristic $0$ or $>2k+3$. For arbitrary values of $a_2$ and of $\beta_2,\beta_3,\beta_4$ (the column $B$), with $a_3:=0$, the columns
\[
C=\frac{1}{ka_2}\bigl[(1+a_2y)^k-1\bigr]
\quad(\text{and } C=y \text{ if } a_2=0),
\qquad
E=E(a_2,\beta_\bullet;y) \text{ as below},
\]
satisfy every key of the depth-two block. Consequently no extra key of pinning type exists, no chart in the coordinates $(a_2,\beta_2,\beta_3,\beta_4)$ is emptied, and the only near-origin restriction on $P$ is the support condition $a_3=0$.
\end{theorem}

\begin{proof}[Proof sketch]
The column-$(k{+}1)$ relation $AC'-kA'C=1$ has the one-dimensional solution family $C=A^k(\kappa+\int_0^yA^{-(k+1)})$, and the valuation $\geq1$ forces $\kappa=0$. The inner extra key is the would-be next series coefficient of this $C$; its vanishing locus contains $\{a_3=0\}$, since with $A=1+a_2y$ the series truncates: $C$ is the stated polynomial of degree $k$, fitting $Q$-column $k$. For column $k+2$, substitute $z=1+a_2y$ and $\beta(z):=B(y(z))$ (a polynomial of $z$-degree $4=2d_2$ with $\beta(1)=0$); the relation $AE'-(k+1)A'E=-(2BC'-kB'C)$ integrates to
\[
E=-\frac{1}{a_2}\,z^{k+1}\!\int_1^z\Bigl[2\beta w^{-3}-\beta'w^{-2}+\beta'w^{-(k+2)}\Bigr]dw .
\]
The first two integrand terms are exact, $2\beta w^{-3}-\beta'w^{-2}=-\frac{d}{dw}(\beta w^{-2})$, and the third integrates to $\sum_m \beta'_m\frac{z^{m-k-1}-1}{m-k-1}$, which is polynomial in $z$ provided no index $m$ equals $k+1$. Since $\deg_z\beta'=2d_2-1=3<k+1$ for $k\geq3$, this always holds: $E$ is a polynomial of degree $\leq k+3\leq 2k+2$, fits $Q$-column $k+1$, and every key of column $k+2$, including both extras, holds identically. The case $a_2=0$ is direct. Every coefficient identity in this argument is additionally machine-checked for $k\leq10$.
\end{proof}

Limitations. The theorem concerns the depth-two block only: deeper Minkowski columns add one leftover key per column, each a residue-type linear condition on the fresh $P$-column; numerical probes suggest these are satisfiable, but we have not proved it. Also, the sharp form of the inner extra key (reduction to $-a_3^k b_{q_0}a_1^{1-k}$, vanishing locus $\{a_3=0\}$ by radicality) is verified by exact computation for $k\leq10$ but not closed for general $k$; the containment used above is proved for all $k$.

\section{A worked example: the $(8,28)$ chart}\label{sec:example}

Proposition 4.3, subcase (2), of \cite{GGHV} asks for pairs with $[P,Q]=x^2$,
\[
N(P)=\mathrm{hull}\{(0,0),(1,0),(8,14),(8,16)\},\qquad
N(Q)=\mathrm{hull}\{(0,0),(2,1),(12,21),(12,24)\},
\]
and the listed vertices attained. This is a reduced strip pair of type $(2,2)$: $d=(1,2)$, $p_0=(1,0)$, $q_0=(2,1)$, strips of lengths $8$ and $12$ columns. The bracket system has $92$ keys; in the repository ordering the vertex key $E_{(3,1)}$ is equation $\#3$,
\[
a_1b_3-1=0,
\]
by uniqueness of the decomposition $(3,1)=(1,0)+(2,1)$ with $\det=1$. As in Theorem \ref{thm:22}, $a_1=a_{(1,0)}$ and $b_3=b_{(2,1)}$.

The pipeline (Section \ref{sec:computation}) runs a level-ordered elimination: after the gap kill $b_{(1,1)}=b_{(1,2)}=0$, it performs $41$ unit-pivot eliminations expressing the $Q$-coefficients of columns $x=12$ down to $x=2$ in $P$-data, then passes to the generic chart. Equation $\#3$ contains only the units $a_1,b_3$ and is never rewritten. The collapse is carried by the key at Minkowski point $(4,7)$ (artifact key $(3,6)$),
\[
E_{(4,7)}=3a_2b_{10}-a_3b_9+2a_5b_6-2a_6b_5+a_{(3,5)}b_{(1,2)}-3a_{(3,6)}b_{(1,1)} .
\]
Ten recorded events touch it, namely the block-local reduction of Theorem \ref{thm:22} plus the gap kill:
\begin{enumerate}
\item[1--2.] $b_{(1,1)}:=0$ and $b_{(1,2)}:=0$ (gap kill);
\item[3--6.] eliminations of $b_{10},b_9,b_8,b_7$ from the keys $(4,6),(4,5),(4,4),(4,3)$;
\item[7--9.] eliminations of $b_6,b_5,b_4$ from the keys $(3,4),(3,3),(3,2)$;
\item[10.] $a_3:=0$ (from the reduced key $(3,5)$).
\end{enumerate}
After event 10 the equation is the single monomial
\[
-\tfrac{1}{5}\,a_2^2\,a_6\,b_3\,a_1^{-2}\;=\;0 .
\]
On the generic chart this coefficient is invertible, so the chart substitutes $b_3:=0$; feeding this into the untouched equation $\#3$ gives $a_1\cdot 0-1=-1$, and the chart is empty. The $-1$ is the right-hand side of the vertex equation; no other arithmetic touches it. After cascade and chart, of the $92$ original equations, $51$ reduce to $0$, $40$ remain polynomial, and equation $\#3$ is the unique constant. The pivot primes are $\{2,3,5,7,11,13\}$ (far-column pivots contribute $7,11,13$; the obstruction itself needs only $\{2,3,5\}$), so the derivation is valid in characteristic $0$ and over $\mathbb{F}_p$ for $p\notin\{2,3,5,7,11,13\}$.

\begin{remark}[Characteristic restrictions]
Mondello's plane Keller pair $P = x + x^2y + x^4 + x^6y^2$, $Q = y + x^5 + x^6y + x^7y^2 + x^8y^3$ over $\mathbb{F}_2$ \cite{Mondello} has Jacobian determinant $1$, collapses three distinct points to one image, and generates a degree $3$ separable extension, so even the separable Jacobian conjecture fails in characteristic two. Its Newton polygons form a strip pair along $(1,1)$ satisfying the vertex normalization. In any characteristic other than two, near-origin eliminations with pivots $2$, $4$, $6$ force the strip corners to vanish and reduce any such pair to an elementary automorphism; in characteristic two those determinants vanish and the strips survive. (Direct bracket check; the example lies in the cell $(k,d_2)=(0,1)$, outside Theorem \ref{thm:22}.)
\end{remark}

In the language of Section \ref{sec:gap}: $(\text{unit})\cdot a_2^2a_6\cdot b_3$ lies in the radical of the system's ideal, so with the vertex equation every solution has $a_2^2a_6=0$, and the variety splits as $\{a_2=0\}\cup\{a_6=0\}$. That split is the case division under which the remaining strata were handed to Gr\"obner engines in the computational discard.

The chain is re-runnable with per-equation provenance and was replayed by adversarial review with independently written geometry, substitution, and census code (no pivot ties; spot checks at random rational points). The two central identities are also checked in Lean~4 with mathlib: the reduction of the key-$(3,6)$ coefficient under the ten substitutions to $(-1/5)\,a_2^2a_6\,\iota_{a_1}^2 b_3$ ($\iota_{a_1}$ the formal inverse of $a_1$), and the collapse of equation $\#3$ to $-1$ at $b_3=0$. The first Lean proof holds already in the free commutative ring and therefore specializes to the localization used here. The formal development covers the identity chain only, not the census or variety preservation of the eliminations.

\section{The log-residue functional: a theorem}\label{sec:functional}

The integral formula in the proof of Theorem \ref{thm:kbig} explains all the coefficients seen so far and, made uniform, proves the general picture. On the binomial locus at any type $(k,d_2)$ ($A=1+a_2y$ after normalizing units, $a_2\neq0$), in the coordinates $z=1+a_2y$ and $\beta(z):=B(y(z))$ (of $z$-degree $2d_2$, $\beta(1)=0$), the only term of the displayed integrand that can obstruct polynomiality of the solved outer column is the index $m=k+1$ of $\beta'=\sum_m\beta'_mw^m$, which integrates to $\beta'_{k+1}\log z$. Since $\deg_z\beta'=2d_2-1$, such a term exists exactly when $k+1\leq 2d_2-1$, and the extra keys see the truncated coefficients of the tail $-\frac{1}{a_2}\beta'_{k+1}z^{k+1}\log z$: the near-origin geometry obstructs when the logarithmic residue $\beta'_{k+1}$ survives. Expanding $\beta$ through $y=(z-1)/a_2$ gives, up to a unit times a power of $a_2$,
\[
\beta'_{k+1}\;\doteq\;R_{k,d_2}\;:=\;\sum_{j=k+2}^{2d_2}(-1)^j\binom{j}{k+2}\,a_2^{\,2d_2-j}\,\beta_j ,
\]
a linear form in the $P$-column-$2$ coefficients $\beta_j=a_{(2,j)}$ with binomial coefficients. This functional organizes every cell; by Theorem \ref{thm:R} below, the computed instances are illustrations, not evidence. At $(2,2)$ the sum has the single term $j=4$, $R_{2,2}=\beta_4=a_6$, and the pinning monomial of Theorem \ref{thm:22} is $a_2^2R_{2,2}$. At $(k,2)$, $k\geq3$, the sum is empty, $R\equiv0$, and Theorem \ref{thm:kbig} confirms that no obstruction exists. At $(2,3)$ and $(2,4)$, $R$ is the hyperplane displayed in Theorems \ref{thm:23full} and \ref{thm:24}. On the binomial locus, in every computed cell one outer extra key survives and equals a unit times a power of $a_2$ times $R_{k,d_2}$, the rest vanishing identically; the verified units $\frac{a_2^3}{21}$ at $(2,3)$, $-\frac{a_2^6}{99}$ at $(3,3)$, $\frac{2a_2^9}{1001}$ at $(4,3)$, $-\frac{a_2^4}{55}$ at $(2,4)$, $-\frac{a_2^8}{364}$ at $(3,4)$, and the identically vanishing outer extras at $(5,3)$, where $k+2>2d_2$ and the sum is empty, are all instances of the closed form in Lemma \ref{lem:outer}.

Everything reduces to one elementary rigidity statement about the weighted derivation identity ($\nu$ is the weight, reserving $w$ for strata; $\operatorname{lc}$ is the leading coefficient).

\begin{theorem}[ODE rigidity, all weights]\label{thm:ode}
Let $F$ have characteristic zero, let $\nu\geq1$ be an integer, and let $A,C\in F[y]$ satisfy $AC'-\nu A'C=c$ for some constant $c\in F$, $c\neq0$. Then $\deg A\leq 1$.
\end{theorem}

\begin{proof}
Both $A$ and $C$ are nonzero, else the left side is $0$; write $\delta:=\deg A$, $\alpha:=\operatorname{lc}(A)$, and suppose $\delta\geq2$.

\emph{Step 1 (leading coefficients).} For nonzero $D\in F[y]$ of degree $d$, the monomial pair $(A_iy^i,D_jy^j)$ contributes $(j-\nu i)A_iD_j\,y^{i+j-1}$ to $AD'-\nu A'D$; the exponent $\delta+d-1$ is reached only by $(i,j)=(\delta,d)$, and nothing higher occurs. So the coefficient of $y^{\delta+d-1}$ is $(d-\nu\delta)\,\alpha\operatorname{lc}(D)$; in characteristic zero, if $d\neq\nu\delta$ then $\deg(AD'-\nu A'D)=\delta+d-1$ exactly.

\emph{Step 2 (degree of $C$).} If $\deg C\neq\nu\delta$, Step 1 gives $\deg(AC'-\nu A'C)=\delta+\deg C-1\geq\delta-1\geq1$, while $\deg c=0$. Hence $\deg C=\nu\delta$.

\emph{Steps 3 and 4 (subtract the kernel).} $A(A^\nu)'-\nu A'A^\nu=0$, so $D:=C-(\operatorname{lc}(C)/\alpha^\nu)A^\nu$ satisfies $AD'-\nu A'D=c$ by linearity, with $\deg D<\nu\delta$. If $D=0$ then $c=0$; if $D\neq0$ then Step 1 gives $\deg(AD'-\nu A'D)\geq\delta-1\geq1$. Either way a contradiction, so $\delta\geq2$ is impossible.
\end{proof}

\begin{remark}[sharpness and prior art]\label{rem:sharp}
The theorem is sharp: $A=1+a_2y$, $C=((1+a_2y)^\nu-1)/(\nu a_2)$ (or $A=1$, $C=y$) solve the identity with $c=1$. The proof is the valuation-at-infinity mechanism of the one-variable constant-term theorem of Duistermaat and van der Kallen \cite{DvdK} made elementary. Characteristic zero enters only through $d-\nu\delta\neq0$: in small characteristic Step 2 genuinely fails, consistent with the characteristic-two example of Section \ref{sec:example}.

The statement has prior art. Dividing the identity by $A^{\nu+1}$ gives $(C/A^\nu)'=c/A^{\nu+1}$, so $A^{-(\nu+1)}$ has a rational primitive; when $A$ has two distinct roots this contradicts Lemma A.7 in the appendix of \.Zo{\l}\k{a}dek \cite{Zoladek}, and a perfect power $A$ dies by evaluating the identity at its multiple root. Lemma A.7 is a non-rationality lemma for Schwarz--Christoffel integrals, phrased there via Darboux functions, which \.Zo{\l}\k{a}dek says was probably known already to Liouville; he applies it inside the plane Jacobian problem via the same ODE, deriving $\phi\psi'-\delta\phi'\psi=\mathrm{const}$, this identity with rational weight $\delta$, from the Jacobian equation in a Newton--Puiseux chart (Lemma 3.9 and equation (3.14) there). The case $\nu=1$, a nonzero constant Wronskian, is Theorem 4 of Hermoso and Alc\'azar \cite{HA}. We proved the theorem independently, by the mechanism above, before locating these references; the proof is retained because the mechanism differs, never leaving $F[y]$, and extends to the positive-characteristic bound $p>(k+1)d_2$ recorded after Theorem \ref{thm:R}.
\end{remark}

The bridge to the block is the solved inner column of Proposition \ref{prop:B}; keep $a_{p_0}=b_{q_0}=1$, so $A=1+a_2y+\dots+a_{d_2+1}y^{d_2}$; the extra keys of Minkowski column $k+1$ are the \emph{inner extras}, those of column $k+2$ the \emph{outer extras} ($d_2-1$ of each).

\begin{lemma}[block to ODE bridge]\label{lem:bridge}
Fix $k,d_2\geq2$ over a field $F$ of characteristic zero and fix a point $a=(a_2,\dots,a_{d_2+1})$. The $d_2-1$ inner extras vanish at $a$ if and only if some $C\in F[y]$ with support in $\{y,\dots,y^{kd_2}\}$ satisfies $AC'-kA'C=1$.
\end{lemma}

\begin{proof}
By \eqref{eq:odes} the inner-column keys are the coefficients of $AC'-kA'C$, $C=\sum_{j=1}^{kd_2}c_jy^j$: the pair $(a_{(1,i)},c_j)$ contributes $(j-ki)\,a_{(1,i)}c_j\,y^{i+j-1}$. The constant coefficient is $c_1$, set to $1$ by the vertex equation. The top coefficient, of $y^{(k+1)d_2-1}$, is $(kd_2-kd_2)\,a_{(1,d_2)}c_{kd_2}=0$ identically: the $w=0$ stratum of Proposition \ref{prop:A}. The coefficients of $y^1,\dots,y^{kd_2-1}$ determine $c_2,\dots,c_{kd_2}$ triangularly (the pivot of $y^n$ is $n+1\neq0$), and the inner extras are the leftover coefficients of $y^{kd_2},\dots,y^{(k+1)d_2-2}$ of the solved series. If the extras vanish at $a$, the solved $C$ satisfies $AC'-kA'C=1$. Conversely a solution with the stated support is unique: two differ by $K$ with $AK'=kA'K$, i.e.\ $(K/A^k)'=0$ in $F(y)$, so $K=\lambda A^k$, and constant terms ($A^k(0)=1$, $K(0)=0$) force $\lambda=0$. Any solution is thus the solved column, and the extras vanish.
\end{proof}

\begin{lemma}[outer column on the binomial locus]\label{lem:outer}
Fix $k,d_2\geq2$, characteristic zero, and let $A=1+a_2y$. Then every outer extra vanishes identically except the lowest (the coefficient of $y^{(k+1)d_2}$), which equals
\[
(-1)^{(k+1)d_2-1}\,\frac{k+2}{\binom{(k+1)d_2}{k+1}}\;a_2^{(k-1)d_2}\,R_{k,d_2}
\]
as a polynomial identity in $(a_2,\beta_\bullet)$, valid at $a_2=0$ as well.
\end{lemma}

\begin{proof}
Write $L(E):=AE'-(k+1)A'E$ and $G:=2BC'-kB'C$, with $C$, $E$ the solved columns of Proposition \ref{prop:B}. On the binomial locus $\deg C=k$, so $\deg G\leq 2d_2+k-1<(k+1)d_2$ (the difference is $(k-1)(d_2-1)>0$), and $\deg A=1$ gives $\deg(L(E)+G)\leq(k+1)d_2$. The coefficients of $y^{d_2},\dots,y^{(k+1)d_2-1}$ are the consumed keys of Proposition \ref{prop:B}, hence zero, and the valuation is at least $d_2$, so $L(E)+G=e\,y^{(k+1)d_2}$ with $e$ the lowest outer extra; the higher outer extras vanish identically. To evaluate $e$, assume $a_2\neq0$ and pair with $\mu(f):=\operatorname{res}_{y=-1/a_2}\bigl(fA^{-(k+2)}dy\bigr)$, the local residue (not the resultant of Section \ref{sec:variants}). Since $L(E)A^{-(k+2)}=(E/A^{k+1})'$ is exact, $\mu(L(E))=0$ and $e\,\mu(y^{(k+1)d_2})=\mu(G)$. Both sides are elementary in $z=1+a_2y$: $\mu(y^n)=(-1)^{n-k-1}\binom{n}{k+1}a_2^{-(n+1)}$, while $G=2\beta z^{k-1}-\beta'z^k+\beta'$ gives $\mu(G)=a_2^{-1}[z^{k+1}]G=(k+2)\,a_2^{-1}[z^{k+2}]\beta$, and $[z^{k+2}]\beta=\sum_j\beta_ja_2^{-j}\binom{j}{k+2}(-1)^{j-k-2}=(-1)^ka_2^{-2d_2}R_{k,d_2}$. Solving for $e$ gives the display. Both sides are polynomials in $(a_2,\beta_\bullet)$, so the identity extends to $a_2=0$.
\end{proof}

\begin{theorem}[log-residue description of the block variety]\label{thm:R}
For every $k\geq2$, $d_2\geq2$, and every reduced strip pair of type $(k,d_2)$ satisfying \textup{(i)}, over a field of characteristic zero, the variety of the extra keys of the depth-two block is
\[
V(\textup{extras})\;=\;\{A \textup{ binomial}\}\;\cap\;\bigl(\{a_2=0\}\cup\{R_{k,d_2}=0\}\bigr),
\]
where ``$A$ binomial'' means $a_{(1,i)}=0$ for $2\leq i\leq d_2$.
\end{theorem}

\begin{proof}
Inner half. If the inner extras vanish, Lemma \ref{lem:bridge} produces a polynomial solution of $AC'-kA'C=1$, and Theorem \ref{thm:ode} with $\nu=k$ forces $\deg A\leq1$: the point is binomial. Conversely on the binomial locus the solution in Remark \ref{rem:sharp} has support in $\{y,\dots,y^k\}$, so the inner extras vanish there. Outer half. On the binomial locus, by Lemma \ref{lem:outer} the outer extras vanish exactly when $a_2^{(k-1)d_2}R_{k,d_2}=0$, i.e.\ when $a_2=0$ or $R_{k,d_2}=0$, since $(k-1)d_2>0$. Intersecting the halves gives the display. When $k+2>2d_2$ the sum defining $R$ is empty, $R\equiv0$, and the statement degenerates to $V(\textup{extras})=\{A\textup{ binomial}\}$ (Theorem \ref{thm:kbig} at $d_2=2$).
\end{proof}

Theorem \ref{thm:R} settles what Section \ref{sec:variants} left open: no point of the depth-two block has all coordinates nonzero when $k,d_2\geq2$ (the forced $a_{(1,2)}=0$); a pinning-type obstruction only at $(2,2)$, where $R$ is a single monomial; no $b$-side obstruction for $d_2\geq3$ ($R$ is then a multi-term linear form, satisfiable with every coordinate nonzero); no residue condition for $k>2d_2-2$. The rigidity halves of Theorems \ref{thm:23} and \ref{thm:24} and the certificates of Theorem \ref{thm:k3} become corollaries, retained as independent corroboration; the cells $(3,4)$, $(4,4)$, $(5,4)$ and all $d_2\geq5$, out of reach of the divisibility argument, are closed (indeed evaluating $AC'-kA'C=1$ at a multiple root gives $0=1$, so perfect powers die with no squarefreeness lemma).

The boundary deserves equal precision. Theorem \ref{thm:R} is a depth-two statement under the strip hypotheses of Definition \ref{def:strip}: Minkowski columns beyond $k+2$ contribute one further residue-type condition each, and their joint behavior remains open. It says nothing at $k=1$ (no gap column, no bridge); nothing for pentagon-shaped or $y$-axis-supported subcases, where the gap kill fails, or for strips of direction $(d_1,d_2)$ with $d_1\geq2$; and it is claimed in characteristic zero only, though the rigidity half holds verbatim for $p>(k+1)d_2$ (a conservative bound on the pivots and weights used).

\section{Application: GGV coverage at $(72,108)$}\label{sec:coverage}

Section 7 of \cite{GGV4} tabulates all $34$ admissible corner families with $\max(\deg P,\deg Q)\leq150$. For each family the type data of a hypothetical strip reduction is pure corner arithmetic: for family corner $A_0=(a,b)$ with $b>a$, the right-hand exponent is $k=\lceil b/a\rceil-2$, and if $b/a$ is not an integer the strip slope is $d_2=a/(\lceil b/a\rceil a-b)$. If $b/a$ is an integer, the $A_0$-corner maps onto the $y$-axis under the reduction, so any subcase keeping that corner has $y$-axis support and is outside Definition \ref{def:strip} (the gap-kill mechanism fails there; see after Lemma \ref{lem:gapkill}). Classifying the $34$ rows:

\begin{center}
\begin{tabular}{lrl}
\toprule
class & rows & basis \\
\midrule
theorem-covered & $1$ & $(8,28)$: strip subcase in the proved $(2,2)$ cell (partial) \\
mechanism-absent, $k=1$ & $19$ & no gap column; condition (ii) fails identically \\
mechanism-absent, $k\geq3$ & $4$ & conditional (see text) \\
discarded & $1$ & $(8,32)$: not a possible pair at all \cite{GGHV} \\
unknown & $9$ & $k=2$ rows, case-by-case reduction not carried out \\
\bottomrule
\end{tabular}
\end{center}

Annotations. The single theorem-covered row is covered only in part: the $(8,28)$ family reduces in \cite{GGHV} to two subcases, and only subcase (2) is strip-shaped; subcase (1) is pentagon-shaped with $y$-axis support, outside the theory, and was handled computationally. For the $19$ rows with $k=1$, the absence of the mechanism is proved (under (i), $k=1$ forces both bottom corners into column $x=1$, so no gap column exists on either side); this does not discard those families, only that this obstruction cannot explain their failure. The $4$ rows with $k\geq3$ all have $b/a$ integral, so if their corners survive reduction the shape is not a strip and the question is moot; if some subcase were strip-shaped, Theorem \ref{thm:R} would cover its depth-two block at every $d_2\geq2$ (with Theorem \ref{thm:kbig} showing that at $d_2=2$ the block imposes nothing beyond support rigidity). Of the $9$ unknown $k=2$ rows, one (the second reduced branch of $(8,28)$ itself) has $d_2=2$ and would land in the proved cell if it strip-reduces; the other eight are $y$-axis flagged.

On current reduction data the proved $(2,2)$ theorem has one client among families of degree $\leq150$: the open $(8,28)$ subcase (2), the case it was built for. No family with degree $\leq150$ is known to need the wide-strip cells $(2,3)$, $(2,4)$, $(k,3)$; Section \ref{sec:variants} is theory completeness and preparation for degrees beyond $150$, not coverage at $150$.

\section{Computational verification}\label{sec:computation}

\noindent\textbf{Pipeline and companion claim.} The campaign behind this note transcribes the Proposition 4.3 polygon data into the full bracket system \eqref{eq:key}, runs the level-ordered unit-pivot elimination and a generic-chart decomposition, and hands each stratum to the Gr\"obner engine msolve \cite{msolve}. The companion claim in the repository is that subcase (2) of Proposition 4.3 has no solutions over any field of characteristic zero, each stratum returning the unit ideal over $\mathbb{Q}$ and at two to three large primes on two independent machines. That claim is conditional on Proposition 4.3 itself; \cite{GGV2,GGV3,GGV4,GGHV} are arXiv-only and unrefereed (versions pinned in the repository, \cite{GGHV} at v1), while \cite{GGV1} is published with one corollary present only in the arXiv version. The theorems of Sections \ref{sec:gap} and \ref{sec:variants} are unconditional statements about strip pairs; only their application to $(72,108)$ inherits the dependency. Theorem \ref{thm:22} replaces the Gr\"obner verdict on the generic chart by a proof and yields the $\{a_2=0\}\cup\{a_6=0\}$ split under which the remaining strata were computed. Subcase (1) is not claimed here.

\medskip
\noindent\textbf{Verification.} Every displayed reduction is checked by re-runnable exact-arithmetic scripts with executable assertions: a full-pipeline replay with per-equation provenance (the ten-event chain of Section \ref{sec:example}), an independent block-local engine that re-derives Propositions \ref{prop:A}, \ref{prop:B} and the reduction tables from corner data alone, and a third layer that solves the column relations \eqref{eq:odes} in exact rational arithmetic, cross-checks them against a lattice-determinant enumeration of every block key, and carries the resultant certificates of Theorems \ref{thm:23}, \ref{thm:24}, \ref{thm:k3}. A further \texttt{mathieu} mode of the residue checker verifies, in exact arithmetic at all computed cells, the identities of Section \ref{sec:functional}: Theorem \ref{thm:ode}'s two computational identities symbolically, the solves of Lemmas \ref{lem:bridge} and \ref{lem:outer}, and the rigidity dichotomy, including every perfect-power instance $(1+uy)^\delta$. Each write-up stage was subjected to adversarial review in a separate session with independently written code. Those reviews confirmed the derivations and produced three corrections incorporated here: the gap condition must be stated side-symmetrically; the chart-free conclusion is radical membership, not ideal membership (Remark \ref{rem:radical}); and numerical probes at $(3,3)$, $(4,3)$, $(5,3)$ had reported solvable systems that are in fact empty (tolerance against multiplicity-three zeros), replaced by the certificates of Theorem \ref{thm:k3}. Numerical verdicts are untrusted throughout; every claim above rests on displayed proofs or exact-arithmetic certificates. Scripts, run logs, review documents, and pinned versions: the archived artifact at \texttt{doi.org/10.5281/zenodo.21894922} (self-contained reproduction bundle: pipeline, certificates, Lean sources, review documents), with the full campaign repository at \texttt{github.com/dcposch/jc72108} (to be made public with this note).

\medskip
\noindent\textbf{AI contribution.} Claude (Fable, Anthropic) performed the eliminations, conjectured the vertex-gap mechanism from the $-1$ collapse, wrote the verification scripts, and drafted this text; independent AI systems (GPT and Grok families) adversarially re-reviewed the mathematics and reimplemented key computations from scratch. The human author directed the work, checked the mathematics, and takes responsibility for the claims. Two of the corrections above came from AI review of AI-produced mathematics.

\small
\begin{thebibliography}{99}

\bibitem{Alpoge} L. Alp\"oge, announcement of an explicit counterexample to the Jacobian conjecture in dimension three (cubic map $\mathbb{C}^3\to\mathbb{C}^3$ with $\det JF\equiv-2$), thread on X, July 20, 2026, \url{https://x.com/__alpoge__/status/2079028340955197566}; detailed write-up announced.

\bibitem{msolve} J. Berthomieu, C. Eder, M. Safey El Din, msolve: a library for solving polynomial systems, Proceedings of ISSAC 2021, ACM, 2021, 51--58.

\bibitem{DvdK} J. J. Duistermaat, W. van der Kallen, Constant terms in powers of a Laurent polynomial, Indag. Math. (N.S.) 9 (1998), 221--231.

\bibitem{GGV1} J. A. Guccione, J. J. Guccione, C. Valqui, On the shape of possible counterexamples to the Jacobian conjecture, J. Algebra 471 (2017), 13--74; arXiv:1401.1784. (The arXiv version contains Corollary 7.12, omitted from the journal version; cf.\ \cite{Heitmann}, Theorem 2.24.)

\bibitem{GGV2} J. A. Guccione, J. J. Guccione, C. Valqui, The two-dimensional Jacobian conjecture and the lower side of the Newton polygon, arXiv:1605.09430.

\bibitem{GGV3} J. A. Guccione, J. J. Guccione, C. Valqui, A system of polynomial equations related to the Jacobian conjecture, arXiv:1406.0886.

\bibitem{GGV4} J. A. Guccione, J. J. Guccione, R. Horruitiner, C. Valqui, Some algorithms related to the Jacobian conjecture, arXiv:1708.07936.

\bibitem{GGHV} J. A. Guccione, J. J. Guccione, R. Horruitiner, C. Valqui, Increasing the degree of a possible counterexample to the Jacobian conjecture from 100 to 108, arXiv:2204.14178v1 (2022).

\bibitem{Heitmann} R. C. Heitmann, On the Jacobian conjecture, J. Pure Appl. Algebra 64 (1990), 35--72.

\bibitem{HA} C. Hermoso, J. G. Alc\'azar, Nonzero constant Wronskians of polynomials and Laurent polynomials, and geometric consequences, arXiv:2410.18867 (2024).

\bibitem{Moh} T.-T. Moh, On the Jacobian conjecture and the configurations of roots, J. Reine Angew. Math. 340 (1983), 140--212.

\bibitem{Shaska} T. Shaska, Graded Keller maps and the Jacobian conjecture, arXiv:2607.20210v2 (2026).

\bibitem{Sigray} I. Sigray, Jacobian trees and their applications, PhD thesis, E\"otv\"os Lor\'and University, Budapest, 2008 (unpublished).

\bibitem{Mondello} G. Mondello, A separable counterexample to the two-dimensional Jacobian conjecture in characteristic two, arXiv:2608.02634 (2026).
\bibitem{Zoladek} H. \.Zo{\l}\k{a}dek, An application of Newton--Puiseux charts to the Jacobian problem, Topology 47 (2008), no.\ 6, 431--469.

\end{thebibliography}

\end{document}
